% Slide 15 — Experimento 1: estudo piloto PB
\begin{frame}{Experimento 1: estudo piloto (Paraíba)}
  \begin{columns}[T,onlytextwidth]
    \column{0.5\textwidth}
      \textbf{Escopo}
      \small
      \begin{itemize}
        \item 15 lojas $\times$ 1 produto (48130)
        \item 12 políticas $\times$ 1 replicação
        \item 180 combinações (política, loja)
        \item Regime \textit{Lumpy} (100\% das lojas)
      \end{itemize}
    \column{0.5\textwidth}
      \textbf{O que validou}
      \small
      \begin{itemize}
        \item Ambiente de simulação e métricas
        \item Protocolo de comparação entre políticas
        \item Comportamentos degenerados (SARSA: FP\,$=$\,0; PPO: FP\,$=$\,0,98)
      \end{itemize}
  \end{columns}
  \vfill
  \begin{alertblock}{Resultado de referência (não generalizável)}
    \small Em volume médio, GA-DQN reduziu CTI em \highlight{48\%} e rupturas
    em 55\% vs.\ $(s,S)$, com BE caindo de 119,4 para 5,0 -- o
    \textbf{ganho máximo observado} neste recorte piloto, não uma média
    universal.
  \end{alertblock}
  \note{
    O Experimento 1 funcionou como piloto: 15 lojas da Paraíba, um único
    produto de alta intermitência, testando o ambiente de simulação e o
    protocolo de comparação com apenas uma replicação. Ele revelou
    comportamentos degenerados em alguns agentes de RL isolados -- o SARSA
    parou de pedir, o PPO pediu em quase todos os ciclos -- o que motivou
    a ampliação para o Experimento 2. O resultado citado de 48% de redução
    de custo do GA-DQN é o ganho máximo observado no estrato de volume
    médio deste piloto; não é, e não deve ser lido como, uma redução média
    universal -- no Experimento 2 esse padrão não se repete da mesma forma.
  }
\end{frame}
